\scanpage{460}\mspage{372}
% archive sheet number: 459
\uline{Corollaire} 4.14.2. \deleted{Avec les notations de} \inserted{Soient $C$ une catégorie} équivalente à une petite catégorie, $\unsure{\mathcal N|N|\mathscr N}$ une \inserted{$\unsure{\mathcal U|U|\mathscr U}$-}catégorie \deleted{accessible}, $\Phi$ et $\Psi$ deux \inserted{petits} ens. de catégories \inserted{$I$}, équivalentes à des petites catégories. On suppose \deleted{\unsure{stables}} $C$ \inserted{et $\mathcal N$} stables par $\varinjlim$ de type $\Phi$, par $\varprojlim$ de type $\Psi$. Soit $\underline{\mathrm{Hom}}_{\Phi}^{\Psi}(C,\mathcal N)$ la sous-catégorie str. pleine de $\underline{\mathrm{Hom}}(C,\mathcal N)$ formée des $f:C\longrightarrow M$ qui commutent aux $\varinjlim$ de type $\Phi$ et aux $\varprojlim$ de type $\Psi$. Alors, si $\mathcal N$ est accessible, il en est de même de $\underline{\mathrm{Hom}}_{\Phi}^{\Psi}(C,\mathcal N)$.

Voici des cas particuliers \deleted{\unsure{plus|bien}}\unsure{remarquables} de ce corollaire :

\uline{Corollaire} 4.14.3. Soient $M,N$ deux catégories accessibles, \inserted{et} $\pi_0$ un cardinal infini, tel que $M$, \deleted{\unsure{et}} $N$ satisfassent à $\unsure{\mathrm L_{\pi_0}|\mathrm L\pi_0}$. \unsure{Soient|Soit} $\underline{\mathrm{Hom}}_{\unsure{!|f|\xi}\pi_0}(M,N)$ la sous-catégorie \unsure{strictmt|strictement} pleine de $\underline{\mathrm{Hom}}(M,N)$ formée des foncteurs $\pi_0$-accessibles. Alors $\underline{\mathrm{Hom}}_{!\pi_0}(M,N)$ est accessible.
